\section{理论指向的实现}
\subsection{先控制标注噪声，再学习策略}
方法将标签质量、表示监督和在线反馈分开处理。
首先由类别定义和少量代表帧生成离散框/点JSON提示，
渲染后人工审核；SAM2\cite{ravi2024sam2}传播种子视频。
通过时间IoU、面积跳变、质心跳变、连通域异常及可用时的双向一致性
定位待重试片段，少量人工复核后训练分视角YOLO提示器，
再扩展到全数据集并迭代修复。
这些诊断分数不是校准概率，也不能发现所有一致性错误。
未来帧仅用于离线标注，不能进入在线反馈。

数据按episode划分。作为全流程未见测试的片段不得参与提示器、
分割器、策略或阈值的拟合；不满足这一条件的旧结果须披露实际已见范围。
质量分数影响离线重试和有效样本选择，并用于几何监督的相对加权。
当前U-Net的类别加权CE与前景Dice不应被写成已默认逐帧质量加权。
分割质量需由独立人工审计支持，而非仅比较模型与其自身伪标签。

\subsection{语义输入与几何监督查询}
为各视角单独训练并冻结轻量U-Net；当前front含背景共7类，
side含背景共5类，保留相应的类别差异。
完整视野缩放到 $480\times270$，
概率 $P^v(c,u)$ 通过固定调色板 $p_c\in[0,1]^3$ 转为
\begin{equation}
 S^v(u)=\sum_c P^v(c,u)p_c.
\end{equation}
这是soft颜色图，不是argmax；多类别到三通道通常不可逆。
每路RGB及对应语义图分别通过同一个ResNet18和投影后进入ACT，
双视角共四路视觉token，单视角两路，
使用原ACT位置编码，不加入camera/modality embedding。
保留RGB是为了避免几何类别丢弃未标注的姿态与接触线索；
共享骨干是实现选择，而不是理论结论。

三个512维可学习query与视觉、关节等token共同进入encoder，
其输出由辅助头预测front几何：
\begin{align}
 G_1&=|M_O|/(W H_I),\nonumber\\
 G_2&=\norm{c_O-c_R}/\sqrt{(W-1)^2+(H_I-1)^2},\\
 G_3&=\bar c_O-\bar e_{T,L}\in\R^2.\nonumber
\end{align}
$c_O,c_R$ 为目标与区域质心；
tool mask在归一化图像平面上做PCA主轴拟合，
选择两个投影极值端点中横坐标更小的 $\bar e_{T,L}$。
该端点是图像约定，不自动等同于真实接触端。

目标不可见时面积零仍监督；关系成员不可见时屏蔽对应关系监督。
令 $v_{bi}$ 表示几何组有效，$q_{bi}$ 为相关类别质量最小值，
$w_{bic}=v_{bi}\min(q_{bi}/0.95,1)$。
当前几何loss按组计算
\begin{equation}
 L_G=\frac13\sum_{i=1}^3
 \frac{\sum_{b,c}w_{bic}\ell_{0.01}(\widehat G_{bic}-\widetilde G_{bic})}
 {\max(1,\sum_{b,c}w_{bic})}.
\end{equation}
$\ell_\beta$ 为Smooth L1；三组维度为1、1、2。
归一化后权重主要改变样本的相对贡献，全组同倍降权可能被分母抵消。
ACT优化 $L_{\rm L1}+\lambda_{\rm KL}L_{\rm KL}+\lambda_G L_G$，
当前 $\lambda_G=1$；几何loss不更新冻结分割器。
动作decoder可读取完整encoder上下文，包括query隐藏状态。
当前没有额外阶段分类、阶段切换或停止头。
区域质心距离也不等价于完整包含率；不能据此声称模型显式判定任务完成。

\subsection{冻结基础策略的历史条件纠偏}
基础策略训练后冻结，残差分支每次只编码新到达的双视角RGB，
完整视野缩放为 $320\times180$，
冻结骨干后池化为每视角 $2\times3\times512$ 特征；
窗口内旧视觉使用缓存，不重复分割或编码。
上下文包括实际follower关节、因果速度估计、上一发送指令、
未来基础动作及相对关节差、有效槽位、计划年龄和观察间隔。

三个实现分支分别为：A使用RGB ACT与最新RGB；
B使用语义ACT，并在快更新时额外提取双视角分割描述；
C使用QToken ACT，将当前生效基础计划的query隐藏状态与最新RGB结合。
C的query不是每次快更新重新计算，不能写为实时几何测量。
A/B/C的基础策略不同，跨分支收益不能直接归因于GRU输入差异。

最近8条因果记录经128维投影与单层hidden128的GRU编码；
每次从窗口起点的零隐藏状态重算，5Hz下跨度约1.4秒。
断流或重置清空历史。快头输出
\begin{equation}
 \delta a_j=\bm\beta\odot\tanh([W_rh+c_r]_j),\qquad
 \widehat a_j=b_j+\delta a_j.
\end{equation}
输出层零初始化，每维限幅由训练集基础误差P95设定，
当前裁剪到 $[0.05,5]$ 原生命令单位，换平台需核对单位。
监督使用leader示范指令、follower观测，不混同此前follower-delta实验。
当前损失为有效槽位上标准差缩放后的Smooth L1
（$\beta=0.1$）及 $10^{-3}$ 的平方残差正则；
AdamW学习率 $10^{-4}$，batch128，最多50轮，验证8轮不改善早停，
以验证缩放MAE选择checkpoint。
理论平方风险另外计算，不从该混合训练loss直接推断成功率。

\subsection{目标槽位对齐与失效处理}
当前30Hz控制，基础动作块60步、每30步重规划；
快更新每6帧一次，目标延迟3帧，输出9步残差。
第6帧观察对应第9至17帧，第12帧观察对应第15至23帧，
相邻修正区间重叠3帧；及时使用首槽的端到端预算为100ms。
所有结果绑定plan ID和目标执行槽位，
只替换尚未执行的对应残差，始终加到原基础动作而非累加旧修正。
计划切换清空旧计划残差，迟到丢弃过期前缀而不移动目标时间。

最大结果观察年龄0.15秒、残差年龄上限0.4秒等规则须以实机部署配置核对。
无有效残差时回退有效基础动作，基础动作也过期时使用已验证的硬件安全处理，
不得无限复用最后指令。实时评测应包含图像获取、预处理、编码、
队列、GRU及发布延迟，而非仅报告神经网络kernel耗时。
